\answer{ex:06:1}{$e=(1-e^{-T_a/\tau_e})e^{-d/\tau_e}$. (أ)~$0.110$ مقابل $4.5\cdot10^{-5}$، أي $\approx2.5\cdot10^3$ مرة. (ب)~$\tau_e^*=T_a/\ln(1+T_a/d)=0.59\,\text{s}$.}
\answer{ex:06:2}{المعدل $\nu_x\nu_y(A_+\tau_+-A_-\tau_-)=0.8A_+$ في الثانية، أي $\approx2900A_+$ في الساعة؛ $\tau_-=\tau_+/0.6=33\ms$.}
\answer{ex:06:3}{$\mathbf e_1$ عند $\mu^2+s_1^2>s_2^2$؛ هنا $1.01>0.25$، فيتعلم العصبون $\mathbf e_1$ مع أن التشتت على امتداد $\mathbf e_2$ أكبر بـ$25$ مرة: القاعدة تجد المتجه الرئيسي لمصفوفة الارتباطات لا التغايرات.}
\answer{ex:06:4}{$V=Re^{-(t_c+L-t)/\tau_\gamma}$ بين المصباح والعصير، فيكون $\dot V=V/\tau_\gamma$؛ ذروة مساحتها $Re^{-L/\tau_\gamma}$: $0.61R$ و$0.78R$.}
\answer{ex:06:5}{نسبة الإشارة إلى الضجيج $1/(N+1)$، وعدد المحاولات $\propto N+1$: $5\cdot10^5$ محاولة $\approx1$ شهر، و$5\cdot10^{11}$ محاولة $\approx8\cdot10^4$ سنة.}
\answer{ex:06:6}{(أ)~$k_2=1/\tau_d$، $k_1=1/\tau_d-1/\tau_\gamma$. (ب)~الخطأ الزائف $-(V/\tau_\gamma)\,\varepsilon/(1+\varepsilon)$، $\varepsilon=\tau_d/\tau_\gamma$، ومنه $\tau_d\le0.105\,\text{s}$.}
\answer{ex:06:7}{$0.46$، $0.90$، $0.99$، $0.95$، و$0.70$ عند $d=3\,\text{s}$. تقع النقاط على منحنى واحد بدلالة نسبة الأثر $F=(1-e^{-T_{\text{cue}}/\tau_e})e^{-d/\tau_e}$: تعلّم عند $F\gtrsim0.1$، واختيار عشوائي عند $F<10^{-2}$.}
\answer{ex:06:8}{$\eta^*=1/\bigl(2\sigma^2(N+2)\bigr)$، في $N+2$ محاولة؛ ومع $m$ عصبونات مرتعشة من $N$ في $N(m+2)/m\ge N+2$: التناثر لا يساعد.}
